% Pursuing Stacks (cotes 134-2 to 134-8): archive cote 134-8, site page 27 (PDF page 28), typescript page 581.
% typescript. Draft compared: ps7.tex:1249-1313.
% AI-transcribed by Claude Opus 5.5 (claude-opus-5-5, adaptive thinking, effort high), Message Batches API, batch ps4 (msgbatch_012Aoc1NuskATPsyBeKSYK7J), 2026-11-12; not yet checked by eye.
\scanpage{134-8/27}
\mspage{581}

such formula, namely the precise relation between $Lf_!$ and $Rf_!$ for
a smooth $f$, was given yesterday \inserted{\deleted{yesterday}} (p.572 (6)).

The day before\hand{,} I was out for trying to get a better understanding
of the $Rf_!$ operation\deleted{\unsure{,}} \hand{(}including the case of non-constant sheaves of
rings for the Modules we work with\hand{)}. This operation still remains unfa-%
miliar to me, very unlike my old friend $Rf_*$ - there is a number of
things about it which are not quite clear \xout{about it} yet in my mind,
even when just taking the functor $f_!$ between Modules, before taking a
left derived functor. For instance, for general rings of operators
$\underline{O}_A$ and $\underline{O}_B$ , when $f$ is a ringed map
\[ \hand{(1)} \qquad f\colon (A,\underline{O}_A)\longrightarrow (B,\underline{O}_B) \quad , \]
it does'nt seem that formation of $f_!$ commutes to restriction of
rings of operators (to the constant rings $\underline{Z}_A$ and $\underline{Z}_B$ say), namely that
for a given $\underline{O}_A$-module $F$, $f_!(A)$ may be interpreted as just $f^{ab}_{o!}(F)$ with
suitable \unsure{operatiions} of $\underline{O}_B$ on the latter, where
\[ \hand{(2)} \qquad f_o : A\longrightarrow B \]
is the map in (Cat) underlying $f$. When reducing to a \inserted{suitable} ``universal'' case,
this may be expressed by saying that for given map $f_o$ and given abelian
sheaf $F$ on $A$, if we define
\[ \margin{\deleted{\illegible{1}}} \qquad \over{\underline{O}}{\underline{O}}_A= \underline{End}_{\underline{Z}}(F) \quad , \quad \underline{O}_B = f_*(\underline{O}_A) \ , \]
there does'nt seem to be a natural operation of $\underline{O}_B$ upon $f^{\hand{ab}}_{\hand{o!}}(F)$; all we
can say, it seems, is that the ring of global sections of $\underline{O}_B$ operates
o\over{n}{n} $f^{ab}_{\hand{o!}}(F)$ . More generally, reverting to the general case of a map $f$
of ringed objects in (Cat), the ring of global sections of $\underline{O}_B$ operates
on $f^{ab}_{o!}(F)$. When $\underline{O}_B$ is a constant sheaf of rings $k_B$ , this implies
of course that $k$ operates on this abelian sheaf on $B$, from which will
follow by an obvious argument that with this structure of a sheaf of
$k$-modules, $f^{ab}_{o!}(F)$ may indeed be identified with $f_!(F)$.

Going over to $Lf_!$ which we would like to express via $Lf^{ab}_{o!}$ , the
situation is worse, as we still have to check (granting \xout{B to be constan}
$\underline{O}_B$ to be constant)that for $F$ a projective $\underline{O}_A$-Module, we got
\[ \hand{(3)} \qquad L_i f^{ab}_{o!}(F)=0 \quad \text{for } i \hand{>} \over{\unsure{0}}{0} \ . \]
This is'nt always true, even when $B$ is the final object in (Cat),
and $A$ has a final object, hence $\underline{O}_A$ is a projective Module over
\over{\illegible{1}}{\unsure{itself}}

\margin{27}
